% GENERATED FILE. Edit canonical lessons, then run npm run generate:books.
% canonical-source-hash: fnv1a64:8551038737f3518a
% canonical-lessons: ZH-C407-jinzhan, ZH-C407-ruanwo, ZH-C407-kongwei, ZH-C407-zhongdianzhan, ZH-C407-dongche, ZH-R407-review-402-407

\chapter{Pulling In}
\label{ch:zh-pulling-in}
\hlchaptermodality{407}

\begin{chapteropening}
\textbf{By the end of this chapter:} \emph{I can say a train is pulling in, name the soft sleeper, find an empty seat and reach the end of the line.}

Complete the last lesson of chapter 407: I can say a train is pulling in, name the soft sleeper, find an empty seat and reach the end of the line.
\end{chapteropening}

\section[\texorpdfstring{\zh{进站} (jìnzhàn)}{进站 (jìnzhàn)}]{\zh{进站} (jìnzhàn) — to pull into the station (of a train or bus)}

\label{lesson:ZH-C407-jinzhan}



\begin{quote}
Before the new one: say the Mandarin for a motor car, then the Mandarin for a luggage rack.
\end{quote}

\subsection*{You'll want to know: \zh{进站}}
\textbf{\zh{进站}} — \emph{jìnzhàn} — \textquotedblleft{}to pull into the station (of a train or bus)\textquotedblright{}.

\begin{quote}
\textbf{\zh{火车进站了现在检票}} \emph{huǒchē jìnzhàn le xiànzài jiǎnpiào} — The train's pulling in: they're checking tickets now.
\end{quote}

\begin{cousinweb}
\textbf{\zh{进站}} (\emph{jìnzhàn}) is built from \zh{进} 'to go forward, to go in' and \zh{站} 'a station', both inherited from Middle Chinese, \emph{tsinH} and \emph{treamH} (Guangyun, in Baxter's transcription).
\end{cousinweb}

\begin{grammarlens}[title={What you've built}]
One more word for getting around.
\end{grammarlens}

\subsection*{Your turn}
\begin{itemize}
\raggedright
  \item \emph{Say it:} \emph{jìnzhàn}
  \item \emph{Say it:} \emph{jìnzhàn}, once more
  \item \emph{Say it:} \emph{huǒchē jìnzhàn le xiànzài jiǎnpiào} — The train's pulling in: they're checking tickets now.
  \item \emph{Recall:} say the Mandarin for a motor car, then the Mandarin for a luggage rack, then say \emph{jìnzhàn} again
\end{itemize}

\begin{tcolorbox}[breakable,colback=teal!4,colframe=teal!35!black,title={Before you move on}]
What does \zh{进站} mean? (\textquotedblleft{}To pull into the station (of a train or bus)\textquotedblright{}.) Say it once more.
\end{tcolorbox}
\section[ruǎnwò]{ruǎnwò — soft sleeper}

\label{lesson:ZH-C407-ruanwo}



\begin{quote}
Before the new one: say the Mandarin for a luggage rack, then the Mandarin for to pull into the station.
\end{quote}

\subsection*{You'll want to know: ruǎnwò}
\emph{ruǎnwò} — \textquotedblleft{}soft sleeper\textquotedblright{}.

\begin{quote}
\emph{ruǎnwò bǐ yìngzuò guì} — A soft sleeper costs more than a hard seat.
\end{quote}

Its characters come later; for now it is a word for the ear.

\begin{cousinweb}
\emph{ruǎnwò} is built from \emph{ruǎn} 'soft' and \emph{wò} 'to lie down', both inherited from Middle Chinese, \emph{nywenX} and \emph{ngwaH} (Guangyun, in Baxter's transcription), which glosses the first 'soft' and the second 'to sleep'.
\end{cousinweb}

\begin{grammarlens}[title={What you've built}]
One more word for getting around.
\end{grammarlens}

\subsection*{Your turn}
\begin{itemize}
\raggedright
  \item \emph{Say it:} \emph{ruǎnwò}
  \item \emph{Say it:} \emph{ruǎnwò}, once more
  \item \emph{Say it:} \emph{ruǎnwò bǐ yìngzuò guì} — A soft sleeper costs more than a hard seat.
  \item \emph{Recall:} say the Mandarin for a luggage rack, then the Mandarin for to pull into the station, then say \emph{ruǎnwò} again
\end{itemize}

\begin{tcolorbox}[breakable,colback=teal!4,colframe=teal!35!black,title={Before you move on}]
What does ruǎnwò mean? (\textquotedblleft{}Soft sleeper\textquotedblright{}.) Say it once more.
\end{tcolorbox}
\section[kòngwèi]{kòngwèi — empty seat}

\label{lesson:ZH-C407-kongwei}



\begin{quote}
Before the new one: say the Mandarin for to pull into the station, then the Mandarin for a soft sleeper.
\end{quote}

\subsection*{You'll want to know: kòngwèi}
\emph{kòngwèi} — \textquotedblleft{}empty seat\textquotedblright{}.

\begin{quote}
\emph{wǒ zài chē lǐ zhǎo bu dào kòngwèi} — I can't find an empty seat on the bus.
\end{quote}

Its characters come later; for now it is a word for the ear.

\begin{cousinweb}
\emph{kòngwèi} is built from \emph{kòng} 'vacant' and \emph{wèi} 'a place', both inherited from Middle Chinese, \emph{khuwngH} and \emph{hwijH} (Guangyun, in Baxter's transcription), which glosses the first 'empty, lacking' and the second as a place in a row.
\end{cousinweb}

\begin{grammarlens}[title={What you've built}]
One more word for getting around.
\end{grammarlens}

\subsection*{Your turn}
\begin{itemize}
\raggedright
  \item \emph{Say it:} \emph{kòngwèi}
  \item \emph{Say it:} \emph{kòngwèi}, once more
  \item \emph{Say it:} \emph{wǒ zài chē lǐ zhǎo bu dào kòngwèi} — I can't find an empty seat on the bus.
  \item \emph{Recall:} say the Mandarin for to pull into the station, then the Mandarin for a soft sleeper, then say \emph{kòngwèi} again
\end{itemize}

\begin{tcolorbox}[breakable,colback=teal!4,colframe=teal!35!black,title={Before you move on}]
What does kòngwèi mean? (\textquotedblleft{}Empty seat\textquotedblright{}.) Say it once more.
\end{tcolorbox}
\section[zhōngdiǎnzhàn]{zhōngdiǎnzhàn — terminus; the end of the line}

\label{lesson:ZH-C407-zhongdianzhan}



\begin{quote}
Before the new one: say the Mandarin for a soft sleeper, then the Mandarin for an empty seat.
\end{quote}

\subsection*{You'll want to know: zhōngdiǎnzhàn}
\emph{zhōngdiǎnzhàn} — \textquotedblleft{}terminus; the end of the line\textquotedblright{}.

\begin{quote}
\emph{wǒmen zài zhōngdiǎnzhàn xià chē} — We get off at the last stop.
\end{quote}

Its characters come later; for now it is a word for the ear.

\begin{cousinweb}
\emph{zhōngdiǎnzhàn} is built from \emph{zhōng} 'the end', \zh{点} 'a point' and \zh{站} 'a station', all inherited from Middle Chinese, \emph{tsyuwng}, \emph{temX} and \emph{treamH} (Guangyun, in Baxter's transcription), which glosses the first 'the utmost, the end'.
\end{cousinweb}

\begin{grammarlens}[title={What you've built}]
One more word for getting around.
\end{grammarlens}

\subsection*{Your turn}
\begin{itemize}
\raggedright
  \item \emph{Say it:} \emph{zhōngdiǎnzhàn}
  \item \emph{Say it:} \emph{zhōngdiǎnzhàn}, once more
  \item \emph{Say it:} \emph{wǒmen zài zhōngdiǎnzhàn xià chē} — We get off at the last stop.
  \item \emph{Recall:} say the Mandarin for a soft sleeper, then the Mandarin for an empty seat, then say \emph{zhōngdiǎnzhàn} again
\end{itemize}

\begin{tcolorbox}[breakable,colback=teal!4,colframe=teal!35!black,title={Before you move on}]
What does zhōngdiǎnzhàn mean? (\textquotedblleft{}Terminus; the end of the line\textquotedblright{}.) Say it once more.
\end{tcolorbox}
\section[dòngchē]{dòngchē — D-class bullet train}

\label{lesson:ZH-C407-dongche}



\begin{quote}
Before the new one: say the Mandarin for an empty seat, then the Mandarin for the end of the line.
\end{quote}

\subsection*{You'll want to know: dòngchē}
\emph{dòngchē} — \textquotedblleft{}D-class bullet train\textquotedblright{}.

\begin{quote}
\emph{zuò dòngchē qù Běijīng yào duōjiǔ} — How long does the bullet train take to Beijing?
\end{quote}

Its characters come later; for now it is a word for the ear.

\begin{cousinweb}
\emph{dòngchē} is built from \emph{dòng} 'to move' and \zh{车} 'vehicle', both inherited from Middle Chinese, \emph{duwngX} and \emph{tsyhae} (Guangyun, in Baxter's transcription), which glosses the first 'to stir, to move'. It is a short form of a longer name, a train whose carriages move under their own power.
\end{cousinweb}

\begin{grammarlens}[title={What you've built}]
One more word for getting around.
\end{grammarlens}

\subsection*{Your turn}
\begin{itemize}
\raggedright
  \item \emph{Say it:} \emph{dòngchē}
  \item \emph{Say it:} \emph{dòngchē}, once more
  \item \emph{Say it:} \emph{zuò dòngchē qù Běijīng yào duōjiǔ} — How long does the bullet train take to Beijing?
  \item \emph{Recall:} say the Mandarin for an empty seat, then the Mandarin for the end of the line, then say \emph{dòngchē} again
\end{itemize}

\begin{tcolorbox}[breakable,colback=teal!4,colframe=teal!35!black,title={Before you move on}]
What does dòngchē mean? (\textquotedblleft{}D-class bullet train\textquotedblright{}.) Say it once more.
\end{tcolorbox}
\section[(dialogue)]{Looking back: from At the Station to Pulling In}

\label{lesson:ZH-R407-review-402-407}



\begin{quote}
Say the Mandarin for a passenger and to wait for a train.
\end{quote}

\subsection*{Your turn}
\begin{itemize}
\raggedright
  \item \emph{Say it:} a passenger, to wait for a train, the high-speed railway, to change trains, a railway carriage
  \item \emph{Say it:} a station waiting room, a passenger on a long journey, on schedule, a train attendant, the head of a station
  \item \emph{Say it:} a ticket office, to check tickets, to buy a ticket on board, a sleeping berth, a hard seat
  \item \emph{Say it:} to go without changing, to go into a place, to go in and out, the last bus or train, to call at a station
  \item \emph{Say it:} crowded, a round trip, to change a booking, a motor car, a luggage rack
  \item \emph{Say it:} to pull into the station, a soft sleeper, an empty seat, the end of the line, a D-class bullet train
  \item \emph{Read it:} \zh{旅} \zh{检} \zh{入} — each one aloud, then a word that holds it
\end{itemize}

\begin{tcolorbox}[breakable,colback=teal!4,colframe=teal!35!black,title={Before you move on}]
A passenger? (\emph{chéngkè}.) To wait for a train? (\textbf{\zh{候车}}, \emph{hòuchē}.) The high-speed railway? (\emph{gāotiě}.) To change trains? (\emph{huànchéng}.) A railway carriage? (\emph{chēxiāng}.) A station waiting room? (\emph{hòuchēshì}.) A passenger on a long journey? (\textbf{\zh{旅客}}, \emph{lǚkè}.) On schedule? (\textbf{\zh{正点}}, \emph{zhèngdiǎn}.) A train attendant? (\emph{lièchēyuán}.) The head of a station? (\textbf{\zh{站长}}, \emph{zhànzhǎng}.) A ticket office? (\emph{shòupiàochù}.) To check tickets? (\textbf{\zh{检票}}, \emph{jiǎnpiào}.) To buy a ticket on board? (\emph{bǔpiào}.) A sleeping berth? (\emph{wòpù}.) A hard seat? (\emph{yìngzuò}.) To go without changing? (\emph{zhídá}.) To go into a place? (\textbf{\zh{进入}}, \emph{jìnrù}.) To go in and out? (\textbf{\zh{出入}}, \emph{chūrù}.) The last bus or train? (\emph{mòbānchē}.) To call at a station? (\emph{tíngkào}.) Crowded? (\emph{yōngjǐ}.) A round trip? (\textbf{\zh{来回}}, \emph{láihuí}.) To change a booking? (\emph{gǎiqiān}.) A motor car? (\emph{qìchē}.) A luggage rack? (\emph{xínglijià}.) To pull into the station? (\textbf{\zh{进站}}, \emph{jìnzhàn}.) A soft sleeper? (\emph{ruǎnwò}.) An empty seat? (\emph{kòngwèi}.) The end of the line? (\emph{zhōngdiǎnzhàn}.) A D-class bullet train? (\emph{dòngchē}.)
\end{tcolorbox}
